% Original: PDF strana 4, donji slajd.
\slajdid{03-s008}
\begin{frame}[label=03-s008]{Izbor komponente signalom CS}
% element: 03-s008:e01
Za izgradnju mreža većeg kapaciteta, sa više ulaza i izlaza, komponente imaju
dodatne signale za izbor komponente: \textbf{CS} (\emph{chip select}).
\par
\begingroup
\def\DEgateScale{.62}\def\DEtapDelta{.45}\def\DEnegSOne{2.1}\def\DEnegSZero{1.2}\def\DEbottomFrame{-.75}
\def\DEdekoderY{0.69}
\def\DEcs{1}\def\DEokvir{1}\def\DEnaziv{Dekoder}\def\DEtip{2/4}\def\DEulaz{CS}
\begin{columns}[c,onlytextwidth]
\begin{column}{.57\textwidth}\centering
% element: 03-s008:e02
% Parametri: DEcs = 0/1; DEokvir = 0/1; DEnaziv, DEtip i DEulaz.
% Font se ne smanjuje skaliranjem kompletnog crteža.
\providecommand{\DEdekoderY}{1}%
\providecommand{\DEgateScale}{.7}%
\providecommand{\DEtapDelta}{.35}%
\providecommand{\DEbottomFrame}{-.47}%
\begin{tikzpicture}[circuit logic US,x=1cm,y=\DEdekoderY cm,DEoznaka,
 inv/.style={not gate US,draw,minimum width=9mm,minimum height=8mm,scale=\DEgateScale},
 kapija/.style={and gate US,draw,minimum width=10mm,minimum height=10mm,scale=\DEgateScale,rotate=-90}]
\providecommand{\DEnegSOne}{2.2}
\providecommand{\DEnegSZero}{1.3}
\foreach \n/\y in {1/2.55,0/1.65}{
 \node[inv] (i\n) at (0,\y) {};
 \node[inv] (j\n) at (1.15,\y) {};
 \draw[DEvod] (i\n.input)--(-1.13,\y) node[left] {S\n};
 \draw[DEvod] (i\n.output)--(j\n.input);
 \coordinate (tap\n) at ($(i\n.output)!.5!(j\n.input)$);
 \draw[DEvod] (j\n.output)--(5.25,\y);
 \draw[DEvod] (tap\n)--++(0,-\DEtapDelta)--(5.25,\y-\DEtapDelta);
 \node[junction] at (tap\n) {};
}
\ifnum\DEcs=1
 \draw[DEvod] (-.95,.95) node[left] {\DEulaz} --(5.25,.95);
\fi
\foreach \n/\x/\sy/\ay in {0/2.05/\DEnegSOne/\DEnegSZero,1/3.05/\DEnegSOne/1.65,2/4.05/2.55/\DEnegSZero,3/5.05/2.55/1.65}{
 \ifnum\DEcs=1
  \node[kapija,logic gate inputs=nnn] (g\n) at (\x,.1) {};
  \draw[DEvod] (g\n.input 3)--(g\n.input 3 |- 0,\sy);
  \node[junction] at (g\n.input 3 |- 0,\sy) {};
  \draw[DEvod] (g\n.input 2)--(g\n.input 2 |- 0,\ay);
  \node[junction] at (g\n.input 2 |- 0,\ay) {};
  \draw[DEvod] (g\n.input 1)--(g\n.input 1 |- 0,.95);
  \node[junction] at (g\n.input 1 |- 0,.95) {};
 \else
  \node[kapija,logic gate inputs=nn] (g\n) at (\x,.1) {};
  \draw[DEvod] (g\n.input 2)--(g\n.input 2 |- 0,\sy);
  \node[junction] at (g\n.input 2 |- 0,\sy) {};
  \draw[DEvod] (g\n.input 1)--(g\n.input 1 |- 0,\ay);
  \node[junction] at (g\n.input 1 |- 0,\ay) {};
 \fi
 \draw[DEvod] (g\n.output)--++(0,-.25) node[below] {Y\n};
}
\ifnum\DEokvir=1
 \draw[DEvod,dashed] (-.87,\DEbottomFrame) rectangle (5.58,3.10);
 \node[align=center] at (.1,.3) {\DEnaziv\\\DEtip};
\fi
\end{tikzpicture}

\end{column}
\begin{column}{.40\textwidth}\centering
% element: 03-s008:e03
% Parametri DEcs, DEnaziv, DEtip, DEulaz.
\begin{tikzpicture}[DEoznaka,x=1cm,y=1cm]
\draw[DEvod] (0,0) rectangle (1.65,2.7);
\node[align=center] at (.82,2.25) {\DEnaziv\\\DEtip};
\foreach \n/\y in {3/1.72,2/1.32,1/.92,0/.52}{
 \draw[DEvod] (1.65,\y)--(1.95,\y);
 \node[anchor=east,inner sep=2pt] at (1.65,\y) {Y\n};
}
\foreach \n/\x in {1/.48,0/1.18}{
 \draw[DEvod] (\x,0)--(\x,-.3);
 \node[anchor=south,inner sep=2pt] at (\x,0) {S\n};
}
\ifnum\DEcs=1
 \draw[DEvod] (-.3,.52)--(0,.52);
 \node[anchor=west,inner sep=2pt] at (0,.52) {\DEulaz};
\fi
\end{tikzpicture}

\end{column}
\end{columns}
\par\endgroup
\par\par\vspace{1mm}
% element: 03-s008:e04
\textbf{Na prikazanom dekoderu aktivna je logička jedinica.}
Kada CS nije aktivan, odnosno \(CS=0\), svi izlazi su na neaktivnom nivou --- logičkoj nuli.
\par\vspace{1mm}
\Rezultat{Komponenta može imati više CS ulaza, povezanih I funkcijom:
\(\textcolor{DEisticanje}{CS=CS_1\cdot CS_2}\).
\textbf{Svi moraju biti aktivni} da bi komponenta bila izabrana.}
\end{frame}
